\documentclass[11pt]{amsart}
\usepackage[margin=1in]{geometry}
\usepackage{amssymb,mathtools,booktabs,mathrsfs}
\usepackage{xurl}
\usepackage[colorlinks=true,linkcolor=blue,citecolor=blue,urlcolor=blue]{hyperref}
\newtheorem{theorem}{Theorem}[section]
\newtheorem{proposition}[theorem]{Proposition}
\newtheorem{lemma}[theorem]{Lemma}
\newtheorem{corollary}[theorem]{Corollary}
\theoremstyle{definition}
\newtheorem{premise}[theorem]{Premise}
\theoremstyle{remark}
\newtheorem{remark}[theorem]{Remark}
\newcommand{\1}{\mathbf 1}
\emergencystretch=4em

\title[A 17-vertex candidate in the Sidorenko construction]{Relaxing the colour-refinement condition in OpenAI's Sidorenko counterexample: a conditional 17-vertex, 30-edge instance}
\author{[AUTHOR NAME]}
\date{October 9, 2026}

\begin{document}
\begin{abstract}
OpenAI's preprint \emph{A counterexample to Sidorenko's conjecture} uses the point--face incidence graph
($35$ vertices, $66$ edges) of a $13$-point triangulated sphere satisfying four conditions (1)--(4).
Condition (4), discrete colour refinement, is used in exactly one place: the pinning lemma. We show
that pinning still works when the stable colouring is not discrete, provided every colour-preserving
homomorphism is an automorphism. The reason is that the coefficient $C$ is invariant under
automorphisms. A $7$-point, $10$-face triangulation satisfies (1)--(3) and this weaker condition. Its
incidence graph $H_7$ has $17$ vertices and $30$ edges. \textbf{Conditionally} on the paper's rank-layer
lemma (and its proof in Sections 4--7) holding for this complex with the same argument, $H_7$ is a
counterexample to Sidorenko's conjecture. This conditional transfer is made at the level of the
\TeX\ argument only and has no machine certificate. It is therefore \textbf{less certain than the
original}, whose $35$-vertex statement is formalised in Lean.
\end{abstract}
\maketitle

\noindent\textbf{Attribution.} The construction and its proof are OpenAI's work (family 161 of
\url{https://github.com/openai/math}, commit \nolinkurl{fd4aeeb2ee4fc729c18d98444fed42fd0529eeeb},
\cite{oai}). This note changes only the choice of the finite complex and one lemma (pinning). It
inherits every other step. Developed with AI assistance (Anthropic Claude, Opus 5.5) under
[AUTHOR NAME]'s direction; independent adversarial reviews were run by separate AI sessions and are not
human peer review.

\medskip
\noindent\fbox{\parbox{0.96\linewidth}{\textbf{Premise and confidence.}
\begin{itemize}
\item[(P)] Lemma 3.1 of \cite{oai} (rank-layer limit) and its proof in Sections 4--7 hold for the complex
$K_7$ below, with the complex-dependent constants re-instantiated (Section~\ref{sec:premise}). This is
\emph{assumed}, not re-proved line by line.
\item The official Lean development certifies only the $13$-point/$22$-face complex (index types are fixed
as \texttt{Fin 13}, \texttt{Fin 22}, \texttt{Fin 33}). Our instance has no machine certificate. Its
reliability is lower than that of the original theorem.
\item Unconditional and exactly checked: (1)--(3) for $K_7$; the colour-preserving homomorphisms of $H_7$;
the coefficient values; the generalised pinning lemma (a short proof, Section~\ref{sec:pinning}).
\end{itemize}}}

\section{Statement}
Let $K_7$ be the complex on points $\{0,\dots,6\}$ with the ten faces in Table~\ref{tab:k7}, and let
$H_7=(I\sqcup J,\mathscr C)$ be its incidence graph: $|I|=7$, $|J|=10$, $|\mathscr C|=30$.
\begin{theorem}\label{thm:main}
Assume (P). There is a finite simple graph $G$ with $t(H_7,G)<p(G)^{30}$, where $p(G)$ is the edge density.
\end{theorem}
The original graph has $35$ vertices and $66$ edges \cite[Thm.~1.1]{oai}.

\begin{table}[h]\small
\begin{tabular}{cccccc}
\toprule
face $j$ & points & opposite faces on $\{i_1i_2\},\{i_1i_3\},\{i_2i_3\}$ & bits\\
\midrule
0&0\,1\,2&1, 2, 6&000\\ 1&0\,1\,6&0, 5, 7&001\\ 2&0\,2\,3&0, 3, 8&001\\ 3&0\,3\,4&2, 4, 8&011\\
4&0\,4\,5&3, 5, 9&111\\ 5&0\,5\,6&4, 1, 7&011\\ 6&1\,2\,5&0, 7, 9&100\\ 7&1\,5\,6&6, 1, 5&110\\
8&2\,3\,4&2, 9, 3&100\\ 9&2\,4\,5&8, 6, 4&110\\
\bottomrule
\end{tabular}
\caption{The complex $K_7$. The edges of $H_7$ are the $30$ pairs (point, face containing it).}
\label{tab:k7}
\end{table}

\section{Conditions (1)--(3) hold; (4) fails}
\begin{proposition}\label{prop:k7}
$K_7$ satisfies conditions (1)--(3) of \cite[Prop.~2.1]{oai}, with $11,13$ replaced by $5,7$:
\begin{enumerate}
\item each of the $|E|=15$ point pairs lies in exactly two faces, and the face-neighbour graph is connected
(Euler characteristic $7-15+10=2$);
\item for each bit position, each class has $5$ faces covering all $7$ points, and it has an exposure order in
which every later face shares a pair with an earlier face and adds exactly one new point:
\[\begin{array}{llll}
\text{bit 1}=0:&0,1,2,3,5\ \ (\text{new } 6,3,4,5)&\text{bit 1}=1:&4,9,6,7,8\ \ (2,1,6,3)\\
\text{bit 2}=0:&0,1,2,6,8\ \ (6,3,5,4)&\text{bit 2}=1:&3,4,5,7,9\ \ (5,6,1,2)\\
\text{bit 3}=0:&0,6,7,9,8\ \ (5,6,4,3)&\text{bit 3}=1:&1,5,4,3,2\ \ (5,4,3,2);
\end{array}\]
\item every $a_e\ge1/3$, and $a_{\{1,6\}}=a_{\{3,4\}}=1$, $a_{\{2,3\}}=a_{\{5,6\}}=2/3$ (all other $a_e=1/3$).
\end{enumerate}
Condition (4) fails. The stable colouring has point classes $\{0\},\{1,4\},\{2,5\},\{3,6\}$ and face classes
$\{012,045\},\{016,034\},\{023,056\},\{125,245\},\{156,234\}$. Moreover $\mathrm{Aut}(H_7)=\{\mathrm{id},\varphi\}$, where
$\varphi$ acts on points by $(1\,4)(2\,5)(3\,6)$ and on faces by $j\mapsto(4,3,5,1,0,2,9,8,7,6)_j$.
\end{proposition}
\begin{proof}
Direct verification from Table~\ref{tab:k7}. Two independent scripts check it exactly
(\texttt{certificate/check\_v7.py}, \texttt{construction/verify3.py}).
\end{proof}

\section{Pinning without discreteness}\label{sec:pinning}
We use the notation of \cite[Section~3]{oai}: base laws $A^{\mathrm{base}}$ with corner signs
$\sigma_\kappa$ (exactly one negative), coefficients $C(a,b)$ of \cite[(3.9)]{oai}, and the scores $\mathcal Q$
of \cite[Lemma~3.4]{oai}.

\begin{lemma}[Pinning with ties]\label{lem:ties}
Let $F=(P\sqcup Q,L)$ be a finite bipartite graph, and run colour refinement from the two part colours until
it is stable (round $R$). Define the scores and $\mathcal Q$ exactly as in the proof of \cite[Lemma~3.4]{oai}.
Then for every part-preserving map $f$ we have $\mathcal Q(f)\le\mathcal Q(\mathrm{id})$, with equality if and only if $f$
preserves the stable colouring.
\end{lemma}
\begin{proof}
The proof of \cite[Lemma~3.4]{oai} shows two things. If $f$ preserves the colours of round $t-1$, then
$F_t(\mathrm{id})-F_t(f)=\sum_u\|h_t(f(u))-h_t(u)\|^2\ge0$, with equality if and only if $f$ also
preserves the round-$t$ colours. And at the first round where a colour is not preserved, the loss outweighs
all later gains. Discreteness was used only to conclude that every $f\ne\mathrm{id}$ has such a round. Without
it, the maps with no such round are exactly the maps preserving the stable colouring, and they tie with the
identity in every round.
\end{proof}

\begin{lemma}\label{lem:hom}
If a part-preserving map $f=(a,b)$ is not a homomorphism $H\to H$, then $C(a,b)=0$.
\end{lemma}
\begin{proof}
Some corner $(i,j)$ is sent to a non-incident type pair $(a(i),b(j))$. There the activation law is
$S=\varnothing$ \cite[Section~3.3]{oai}, so $A_{a(i)b(j)}(\alpha,\beta)=0$ unless $\alpha=\beta=\varnothing$. Every term in
\cite[(3.9)]{oai} has all labels nonempty.
\end{proof}

\begin{lemma}[Automorphism invariance]\label{lem:aut}
For every automorphism $\varphi$ of $H$ (preserving the parts),
$C(\varphi)=\bigl(\prod_\kappa\sigma_\kappa\bigr)(\theta/2)^{|\mathscr C|}=C(\mathrm{id})$.
\end{lemma}
\begin{proof}
An automorphism permutes corners, pairs and faces compatibly, so $\varphi(U_e)=U_{\varphi(e)}$,
$\varphi(M_\kappa)=M_{\varphi(\kappa)}$ and $\varphi(\Gamma_\kappa)=\Gamma_{\varphi(\kappa)}$. Relabelling the coordinate set by $\varphi^{-1}$
gives $A^{\mathrm{base},\sigma}_{\varphi(\kappa)}(\alpha,\beta)=A^{\mathrm{base},\sigma\circ\varphi}_{\kappa}(\varphi^{-1}\alpha,\varphi^{-1}\beta)$, and hence
$C_\sigma(\varphi)=C_{\sigma\circ\varphi}(\mathrm{id})$. The support-forcing Lemma~3.3 of \cite{oai} holds for every sign vector
(its proof uses only (1) and the structure of $M_\kappa,\Gamma_\kappa$), with value $\prod_\kappa(\sigma\circ\varphi)_\kappa\,(\theta/2)^{|\mathscr C|}$. The
product of the signs is unchanged by the permutation.
\end{proof}

\begin{proposition}\label{prop:maxhom}
The colour-preserving part-preserving maps of $H_7$ number $2^6\cdot2^{10}=65536$. Exactly two of them are homomorphisms,
namely $\mathrm{id}$ and $\varphi$.
\end{proposition}
\begin{proof}
Exhaustive enumeration. Three implementations agree: \texttt{construction/verify3.py} (all $65536$ maps),
\texttt{construction/pinexact.py} (variable elimination over all part-preserving maps and over all
homomorphisms), and \texttt{certificate/check\_v7.py} (written independently from the \TeX: $32033$ point maps
extend to homomorphisms; exactly two maximise $\mathcal Q$, with exact $C=-(\theta/2)^{30}$ each).
\end{proof}

\begin{proof}[Proof of Theorem~\ref{thm:main}]
In \cite[(3.21)]{oai}, replace $13,22,66$ by $7,10,30$. By Lemmas~\ref{lem:ties}--\ref{lem:aut} and
Proposition~\ref{prop:maxhom}, the maximal-score part of the sum is $C(\mathrm{id})+C(\varphi)=-2(\theta/2)^{30}$. Every other
map has score at most $\mathcal Q_*-\delta_*$ with $\delta_*>0$. Choosing $\rho$ with $B_*e^{-\rho\delta_*}<2(\theta/2)^{30}$ gives $t_0<0$. Dilution gives
$t(W_\varepsilon)=1+\varepsilon^{7}t_0+o(1)$ and $t(W_\varepsilon^\top)=1+\varepsilon^{10}t_1+o(1)$, so the bracket of
\cite[(3.25)]{oai} becomes $t_0+\varepsilon^3t_1+\varepsilon^{10}t_0t_1$. The crossed product and the finite-host transfer
\cite[Prop.~3.5, Lemma~3.6]{oai} apply unchanged. All of these steps use Lemma~3.1, hence (P).
\end{proof}

\section{What the premise covers}\label{sec:premise}
In \cite{oai}, condition (4) is cited only in the application of the pinning lemma (\texttt{kernel.tex} l.~455).
Sections 4--7 use (1)--(3) and numbers derived from the complex. Re-instantiated for $K_7$, these are:
$33\to15$ pairs; $22\to10$ faces in the slack bound \cite[\S7]{oai} ($22C_*\to10C_*$, so
$D>M_*+20C_*$ suffices there); point degrees $(6,4,5,3,4,5,3)$, so $\sum_i\binom{2\deg i}{2}=242$ (was $618$); and the
profile exponent bound $m^2-4m+242\le407$ for $m\le15$ (was $1575$). A separate review session checked these
substitutions. We did \emph{not} re-derive Sections 4--7 for $K_7$; this is the content of (P).

\section{Controls: smaller complexes}
Removing (4), every sphere triangulation with $V\le6$ satisfying (2)--(3) fails for a definite reason.
\begin{itemize}
\item $V=4,5$ and the octahedron ($V=6$, \#2): $H$ is already a Sidorenko graph. A reflection (an involution of $H$
splitting the faces into two halves glued along fixed points) and the Cauchy--Schwarz inequality reduce $H$ to a half with
a part of size $\le4$ (or a dominating vertex), which is Sidorenko \cite{sidorenko1991,cfs2010}
(\texttt{construction/survey.py}). The tetrahedron's $H$ has a part of size $4$.
\item $V=6$, \#1 (faces $012,015,023,034,045,124,145,234$): the pinning step fails. Eight colour-preserving
\emph{folding} homomorphisms have $C=+\theta^{22}/2^{24}$, which dominates $C(\mathrm{aut})=-\theta^{24}/2^{24}$ for small $\theta$. They form pairs with
$\mathrm{vec}(f)+\mathrm{vec}(f')=2\,\mathrm{vec}(\mathrm{id})$ (counts of point, face and corner images), so by AM--GM no choice of priors and thinning weights the
identity above both (\texttt{construction/fold1.py}). Whether this $H$ is Sidorenko is not decided here.
\end{itemize}
At $V=7$ the other two complexes behave differently. \#4 is Sidorenko by reflection. For \#5, generalised
pinning is inconclusive ($270$ maximising homomorphisms against $|\mathrm{Aut}|=6$). So $K_7$ (\#3) is the smallest instance
in the relaxed template for which this argument goes through. The same argument also works for $V=8$ (\#6) and $V=9$ (\#10).

\section{$H_7$ is not in a known Sidorenko class}
\begin{table}[h]\small
\begin{tabular}{p{0.42\linewidth}p{0.5\linewidth}}
\toprule
class (positive theorem) & why $H_7$ is not covered\\
\midrule
trees, even cycles, complete bipartite; one part of size $\le4$ \cite{sidorenko1991} & parts $7,10$; cycles; min degree $3$\\
vertex complete to the other part \cite{cfs2010} & max degrees $6<10$ and $3<7$\\
weakly norming (needs regular parts, \cite{hatami2010}) & point degrees $3,3,4,4,5,5,6$\\
Coregliano's degree-count criterion \cite{coregliano2024} & $10<\binom73=35$; reverse orientation also fails\\
blow-ups \cite{conlonlee2021} & no twin vertices\\
Cartesian products & $|V(H_7)|=17$ is prime\\
subdivisions / even-path substitutions & min degree $3$\\
reflection (symmetric gluing) & $\mathrm{Aut}=\mathbb Z/2$; $\varphi$ fixes only point $0$, and $H_7-0$ is connected\\
strongly tree-decomposable \cite{cklee2018}, reflection trees, tree-arrangeable & first-step leaf test: $56$ candidate leaf sets, none passes (\texttt{classes/stdcheck.py})\\
Szegedy's ``thick'' class & \textbf{not decided}\\
\bottomrule
\end{tabular}
\caption{Location of $H_7$; invariants in \texttt{classes/invariants.py}. Conlon--Lee's further sufficient
conditions (divisibility, $|B|\ge\binom{|A|}{r}$) were checked by the review session to fail as well.}
\end{table}
Also, a numerical search over $69$ step-kernel starts (\texttt{wsearch}) always converged to the constant kernel.
The review session reports the same behaviour for the original graph: the counterexample kernel is a high-dimensional rank-layer kernel, not a small step function.

\section{Relation to other work}
The public repository \texttt{rohanarun/sidorenko-template-refinements} (2026-10-09, commit \texttt{4e1b9b6})
keeps condition (4) and proves that, within that template, the $13$-point/$35$-vertex complex is minimal.
For $V\le12$ they give exact cut-obstruction certificates. Their crossing-sum identity
$\sum_{e\ni v}a_e=2$ forces point degree $\le6$, and $K_7$ attains this ($a_e=1/3$ on all six pairs at point $0$).
Relaxing (4) is exactly what takes us past their template minimum. Appendix~\ref{app:lp} is an
independent recomputation of their $V\le10$ range.

\appendix
\section{Independent LP certificates with (4) kept}\label{app:lp}
\begin{lemma}\label{lem:ham}
For a sphere triangulation, the admissible bit splits (both classes have spanning exposure orders) correspond
bijectively, up to swapping the classes, to Hamiltonian cycles of the point graph. The cycle is the set of pairs
separated by the split.
\end{lemma}
\begin{proof}
An exposure order glues each new triangle along a boundary edge of a disc and adds a new vertex. Each class is
therefore a triangulated disc with all $V$ points on its boundary, and the common boundary is a Hamiltonian
cycle. Conversely, a Hamiltonian cycle bounds two discs with no interior vertices (Jordan). Reversing an
ear decomposition of each disc gives exposure orders.
\end{proof}
Fractional (2)+(3) asks for a probability vector $w$ on Hamiltonian cycles covering every pair with mass $\ge1/3$. A
rational certificate $y\ge0$ with $\sum_ey_e=1$ and $y(C)<1/3$ for every Hamiltonian cycle $C$ rules this out.
\texttt{appendix\_lp/r6\_verify.py} enumerates all triangulations with $V\le10$ (plantri 5.8) and finds the discrete ones:
$0,0,0,0,2,16,137$ for $V=4,\dots,10$. It certifies all $155$ in \texttt{fractions.Fraction} against every Hamiltonian cycle; the
worst certified bounds are $1/5,2/7,2/7$ at $V=8,9,10$. As a positive control, the original complex is feasible with
$t^*=1/3$. Restricting to point degree $\le6$ gives $0,1,4$ at $V=8,9,10$, the same as in
\texttt{rohanarun}. An earlier run (\texttt{primary/}, plantri 5.5) covered $V\le12$ ($7881$ discrete complexes, all
infeasible with exact certificates). Its $V=11,12$ part has not been re-run independently.

\begin{thebibliography}{9}
\bibitem{oai} OpenAI, \emph{A counterexample to Sidorenko's conjecture}, OpenAI Math Release preprint, September 23, 2026,
\url{https://github.com/openai/math/tree/main/preprints/A-counterexample-to-Sidorenkos-conjecture-September-23-2026}.
\bibitem{sidorenko1991} A.~Sidorenko, Inequalities for functionals generated by bipartite graphs, Diskret.\ Mat.\ 3 (1991); and A correlation inequality for bipartite graphs, Graphs Combin.\ 9 (1993).
\bibitem{cfs2010} D.~Conlon, J.~Fox, B.~Sudakov, An approximate version of Sidorenko's conjecture, GAFA 20 (2010).
\bibitem{hatami2010} H.~Hatami, Graph norms and Sidorenko's conjecture, Israel J.\ Math.\ 175 (2010).
\bibitem{cklee2018} D.~Conlon, J.~H.~Kim, C.~Lee, J.~Lee, Some advances on Sidorenko's conjecture, J.\ London Math.\ Soc.\ 98 (2018).
\bibitem{conlonlee2021} D.~Conlon, J.~Lee, Sidorenko's conjecture for blow-ups, Discrete Anal.\ (2021).
\bibitem{coregliano2024} L.~N.~Coregliano, as cited in \cite[\S1]{oai} (Theorem~3.6 there).
\end{thebibliography}
\end{document}
